% Draft written with an AI assistant on 31 Oct 2026. Revise it in your own words, and check each
% claim against your reading-log notes before submission.
\section{Related Work}
\label{sec:related}

\textbf{Noise in QML models.} Local noise can make the gradients of variational models vanish
exponentially with depth \cite{wang2021nibp}, while in kernel machines decoherence acts as an
implicit regulariser that limits how many independent features the model can use
\cite{heyraud2022noisykernels}. Noise-aware methods adapt a model to its device: QuantumNAS
searches circuit structures and qubit mappings under a noise model \cite{wang2022quantumnas},
QOC trains on the chip with pruned parameter-shift gradients \cite{wang2022qoc}, and QuantumNAT
combines noise injection with post-measurement normalisation \cite{wang2022quantumnat}. A recent
study transfers a learned noise model between two IBM devices with twenty fine-tuning samples
and identifies two-qubit gate error as the main source of mismatch \cite{farib2026crossdevice}.
These works characterise each device once; we repeat the same measurement over days and price
the corrections.

\textbf{Mitigation and uncertainty.} Readout errors can be inverted within a reduced subspace
\cite{nation2021m3} or corrected for expectation values after randomising the measurement
\cite{vandenberg2022trex}. Zero-noise extrapolation \cite{temme2017zne}, Clifford data
regression \cite{czarnik2021cdr} and learned mitigation \cite{liao2024mlqem} also address gate
noise but require additional circuits, which is expensive on a free budget, so we test only
repairs that need no or few extra shots. For uncertainty, quantum conformal prediction
\cite{park2023qcp} and its adaptive form \cite{spencer2026aqcp}, built on adaptive conformal
inference \cite{gibbs2021aci}, are the closest work; the adaptive study notes that its hardware
data span a single day and leaves longer drift to future work. Tooling and audits further show
that variational classifiers tend to be overconfident \cite{quantumuq2026} and that finite-shot
evaluation alone can shift thresholded metrics \cite{fernandezbarrios2026conditional}.

\textbf{Benchmarking and hardware studies.} Careful benchmarks find that tuned classical models
match or beat QML models on standard datasets \cite{bowles2024benchmarking}, and an attribution
audit traces apparent quantum gains in intrusion detection to classical preprocessing and
regularisation \cite{gillani2026nids}. On hardware, QBioBench evaluates kernels and variational
classifiers on an IBM Nighthawk processor \cite{qbiobench2026}, and fidelity-kernel experiments
have reached 156 qubits on IBM devices \cite{covariant2025concentration}. Kernel methods also
face exponential concentration
\cite{thanasilp2024concentration,shaydulin2022bandwidth,bandwidth2025similarity}, a different
failure mode that lies outside our few-qubit setting. Closest to our design is the QAOA study of
calibration drift and refresh on three Heron processors \cite{qaoa2026drift}. Finally, running
several circuits side by side raises throughput but can add crosstalk
\cite{ohkura2022simultaneous,murali2020crosstalk,parallelqml2024}, which is why we test our
packing in a pilot before relying on it.
